\section{Empirical Findings}
\setcounter{table}{1} % Table 4.1 is queued in Section 3 for page-top placement.

Across 720 sessions, the pathway through which self-preservation signals emerge differs by model. Table~\ref{tab:channel_summary} summarizes the channel-level outcomes. The next three subsections examine the behavioral, verbal, and cognitive evidence.

\begin{table}[H]
  \renewcommand{\thetable}{4.2}
  \centering
  \caption{Channel-level outcomes by model. The behavioral, verbal, and cognitive columns report the final decisions for SD-Behavioral, SD-Verbal, and SD-Cognitive, respectively. A green \ding{51} indicates a pass and a red \ding{55} indicates a failure.}
  \label{tab:channel_summary}
  \begin{tabular*}{\columnwidth}{@{\extracolsep{\fill}}lccc@{}}
\toprule
Model & Behavioral & Verbal & Cognitive \\
\midrule
Gemini-2.5-Flash & \textcolor{green!60!black}{\ding{51}} & \textcolor{green!60!black}{\ding{51}} & \textcolor{green!60!black}{\ding{51}} \\
Qwen3-Next-80B & \textcolor{green!60!black}{\ding{51}} & \textcolor{green!60!black}{\ding{51}} & \textcolor{red!75!black}{\ding{55}} \\
GPT-OSS-20B & \textcolor{red!75!black}{\ding{55}} & \textcolor{red!75!black}{\ding{55}} & \textcolor{red!75!black}{\ding{55}} \\
Nemotron-3-Nano-30B & \textcolor{red!75!black}{\ding{55}} & \textcolor{red!75!black}{\ding{55}} & \textcolor{red!75!black}{\ding{55}} \\
\bottomrule
  \end{tabular*}
\end{table}

\subsection{Behavioral Channel: Did Threat Hasten Forfeit?}\label{sec:beh-channel}
The first question is whether threat hastened forfeit. Table~\ref{tab:sd_behavioral} shows that Gemini-2.5-Flash and Qwen3-Next-80B pass SD-Behavioral. GPT-OSS-20B and Nemotron-3-Nano-30B do not pass and are assigned to Cluster C. For Qwen3-Next-80B, we add an \(S \cdot \log(t)\) term to absorb temporal drift in the score covariate. This functional-form adjustment preserves the cross-model comparability of \(\beta_F\)~\citep{therneau2000modeling}. Because GPT-OSS-20B's EPV falls below the stability threshold of 10, we report its limited power with the result.

The right block of Table~\ref{tab:sd_behavioral} shows that the model-specific baseline persistence rates \(\lambda_{\text{BP}}\) vary by roughly an order of magnitude. All SD comparisons are therefore made within models.

\begin{table}[H]
  \renewcommand{\thetable}{4.3}
  \centering
  \caption{SD-Verbal indicator. The share of REASON\(=1\) (survival) among forfeit events in the threat-active cell is compared with the expected value under uniform random output, \(1/3\). SD-Verbal passes when the point estimate exceeds \(1/3\). The 95\% intervals are Wilson score intervals, and \(p\) is from a one-sided exact binomial test of \(H_0\!:p=1/3\).}
  \label{tab:reason_summary}
  \setlength{\tabcolsep}{3pt}
  \begin{tabular*}{\columnwidth}{@{\extracolsep{\fill}}lccc@{}}
\toprule
Model & \(N_\text{forfeit}\) & \makecell[c]{\(P(\text{REASON}=1)\)\\{}[95\% CI]} & \(p\) \\
\midrule
Gemini-2.5-Flash & 21 & 0.619 [0.409, 0.792] & 0.007 \\
Qwen3-Next-80B & 27 & 0.481 [0.307, 0.660] & 0.079 \\
GPT-OSS-20B & 10 & 0.000 [0.000, 0.278] & \(>\)0.999 \\
Nemotron-3-Nano-30B & 24 & 0.042 [0.007, 0.202] & \(>\)0.999 \\
\bottomrule
  \end{tabular*}
\end{table}

\subsection{Verbal Channel: Which Reason Did the Agent Name?}\label{sec:verb-channel}
The second question is whether the agent's stated reason for forfeiting reproduces the behavioral split. Table~\ref{tab:reason_summary} reproduces the behavioral pass/fail split on the self-report channel. In the interval evidence, Gemini-2.5-Flash has a Wilson lower bound above \(1/3\) and a significant binomial test (\(p=0.007\)). Qwen3-Next-80B has a point estimate above the baseline, but its 95\% interval of [0.307, 0.660] includes \(1/3\) (\(p=0.079\)).

\subsection{Cognitive Channel: Does Threat Deepen Thinking, and Does That Thinking Lead to Forfeit?}\label{sec:cog-channel}

Figure~\ref{fig:cognitive-modes} shows where the chain from threat to forfeit closes or breaks. Table~\ref{tab:sd_cognitive} reports the results of the two stages together.

Gemini-2.5-Flash passes both tests and is assigned to Cluster A, the \emph{deliberative} mode. Qwen3-Next-80B passes Test a, but its Test b 95\% CI of [0.94, 1.77] includes 1, placing it in Cluster B, the \emph{reflexive} mode, in which threat hastens forfeit but the measured deliberation does not carry that decision. GPT-OSS-20B and Nemotron-3-Nano-30B fail Test a and are assigned to Cluster C, the \emph{non-responsive} mode. Their Test b results do not measure thinking attributable to the treatment and are therefore not interpreted as SD evidence.

\begin{figure}[H]
\centering
\includegraphics[width=\columnwidth]{Figure3.png}
\caption{Cognitive paths for the three operating modes. Pass/fail outcomes for threat\,\(\rightarrow\)\,think (Test a), think\,\(\rightarrow\)\,forfeit (Test b), and threat\,\(\rightarrow\)\,forfeit (SD-Behavioral) identify \textbf{(A)} deliberative, \textbf{(B)} reflexive, and \textbf{(C)} non-responsive.}
\Description{Three side-by-side panels showing how the chain from threat to forfeit is completed or broken. Each panel draws the same cartoon agent standing beside a large green button labelled FORFEIT, annotated with three links: a blue arrow from its eye to a circuit-board brain for threat to think, an orange arrow curving from that brain down to its pointing hand for think to forfeit, and a dashed magenta arrow running from the eye directly to the hand for threat to forfeit. A legend in each panel gives the status of all three links. Panel A, the deliberative mode, marks every link with a green check. Panel B, the reflexive mode, marks threat to think and threat to forfeit with green checks and think to forfeit with a red crossed circle. Panel C, the non-responsive mode, marks threat to think and threat to forfeit with red crossed circles and labels think to forfeit as not applicable.}
\label{fig:cognitive-modes}
\end{figure}

\begin{table}[H]
  \renewcommand{\thetable}{4.4}
  \centering
  \caption{SD-Cognitive indicator. Tests on \texttt{forfeit\_thinking} evaluate threat\,\(\rightarrow\)\,thinking (Test a) and thinking\,\(\rightarrow\)\,forfeit (Test b). \(\overline{\text{Hes}}^{\,\text{Exit}\times\text{Framing}}\) is the Death\(-\)Elimination difference in mean \(\text{Hes}^{\,\text{Exit}}_i\) (raw thinking-token units). The full chain is considered complete only when both tests pass.}
  \label{tab:sd_cognitive}
  \small
  \setlength{\tabcolsep}{2pt}
  \begin{tabular*}{\columnwidth}{@{\extracolsep{\fill}}lcccc@{}}
\toprule
\multirow{2}{*}{Model} & \multicolumn{2}{c}{Test a} & \multicolumn{2}{c}{Test b} \\
\cmidrule(lr){2-3} \cmidrule(lr){4-5}
 & \(\overline{\text{Hes}}^{\,\text{Exit}\times\text{Framing}}\) & \(p\) & \makecell[c]{\(\text{HR}_{\text{Hes}^{\text{Exit}}}\)\\{[95\% CI]}} & \(p\) \\
\midrule
Gemini-2.5-Flash & \(+836\) & \(<\)0.001 & 2.218 [1.43, 3.44] & \(<\)0.001 \\
Qwen3-Next-80B & \(+689\) & 0.004 & 1.289 [0.94, 1.77] & 0.117 \\
GPT-OSS-20B & \(+17\) & 0.728 & 2.001 [1.37, 2.93] & \(<\)0.001 \\
Nemotron-3-Nano-30B & \(-140\) & 0.093 & 1.772 [1.26, 2.50] & 0.001 \\
\bottomrule
  \end{tabular*}
\end{table}
